% Generated draft: requires mathematical and source review.
\documentclass[11pt]{article}
\usepackage[margin=25mm]{geometry}
\usepackage{fontspec}
\setmainfont{Latin Modern Roman}
\usepackage{amsmath,amssymb,mathtools,unicode-math}
\setmathfont{Latin Modern Math}
\usepackage{enumitem,needspace,xurl,hyperref}
\setlength{\parindent}{0pt}
\setlength{\parskip}{5pt}
\setlength{\emergencystretch}{3em}
\newcommand{\N}{\mathbb{N}}
\newcommand{\Z}{\mathbb{Z}}
\newcommand{\Q}{\mathbb{Q}}
\newcommand{\R}{\mathbb{R}}
\newcommand{\C}{\mathbb{C}}
\newcommand{\F}{\mathbb{F}}
\newcommand{\A}{\mathbb{A}}
\newcommand{\PP}{\mathbb{P}}
\newcommand{\E}{\mathbb{E}}
\newcommand{\eps}{\varepsilon}
\newcommand{\dd}{\,\mathrm{d}}
\begin{document}
\section*{UnsolvedMath ID: 800930}
\textbf{Source problem code:} catalogue entry
\subsection{Definitions and mathematical statement}
Let \(M\) be a compact manifold with a Riemannian metric \(g\), or with a Finsler metric in the Finsler version, and let \(G\) be a finitely generated abelian group. Write \(\Lambda M\) for its free loop space and \(L_g(\gamma)\) for the length of a loop \(\gamma\). For \(a\geq 0\), set
\[
\Lambda_a M=\{\gamma\in\Lambda M:L_g(\gamma)\leq a\},
\]
and let \(i_a:\Lambda_a M\hookrightarrow\Lambda M\) be the inclusion. For \(Y\in H_*(\Lambda M;G)\), define its critical level by
\[
\operatorname{cr}_g(Y)=\inf\bigl\{a\geq 0:Y\in\operatorname{im}\bigl((i_a)_*:H_*(\Lambda_a M;G)\longrightarrow H_*(\Lambda M;G)\bigr)\bigr\}.
\]
In particular, this definition gives \(\operatorname{cr}_g(0)=0\). For a nonzero class \(Y\), the supplied source also gives the minimax description
\[
\operatorname{cr}_g(Y)=\inf_{\substack{c\text{ a singular cycle}\\ [c]=Y}}\ \sup_{\gamma\in |c|}L_g(\gamma),
\]
where \(|c|\) is the union of the images of the singular simplices occurring with nonzero coefficients in \(c\). The source does not specify the regularity or topology of \(\Lambda M\); a precise analytic model supporting these length sublevels and the stated minimax description remains to be fixed in review.

For \(M=S^n\), with \(n\) retained as an unspecified dimension parameter, the question attributed to Hingston is whether there exist a metric \(g\), a class \(0\neq X\in H_*(\Lambda S^n;\Z)\), and a positive integer \(m\) such that
\[
0<\operatorname{cr}_g(mX)<\operatorname{cr}_g(X).
\]
Here \(mX\) is the integer multiple of the homology class \(X\), rather than an operation that iterates loops. The positive lower bound necessarily excludes \(mX=0\). The supplied record does not specify the intended range or quantifier of \(n\), or whether the final existence question is intended separately for Riemannian and Finsler metrics. These qualifications remain unresolved.
\subsection{Short English statement}
For a sphere of the intended dimension, can multiplying an integral homology class of its free loop space by a positive integer strictly lower its minimax length while keeping that length positive? The dimension quantifier, the precise loop-space model, and the intended Riemannian versus Finsler scope require clarification.
\subsection{Sources}
\begin{itemize}
\item Supplied local UnsolvedMath registry record: permanent ID \(800930\), source code \texttt{catalogue entry}, collection \emph{AMR Open Problem Lists}. The record labels the problem open; this is a supplied status label, not an independently verified assessment.
\item Burns and Matveev, \emph{Open Problems and Questions about Geodesics} (2019, as identified in the supplied record), Question 10. Supplied source URL: \url{https://arxiv.org/abs/1308.5417}. The record attributes the question/context to Hingston.
\item The supplied provenance reports access on 2026-10-16 and extraction by \texttt{arxiv\_latex}. These provenance details are reproduced from the record; no independent source visit or literature verification is asserted.
\end{itemize}
\end{document}
